\begin{abstract}
Breadth-first search (BFS) is a basic graph algorithm. The fastest parallel BFS implementations for shared memory are direction-optimizing: they expand a search level by level, each level top-down or bottom-up. They are also blocking. Every level or round ends at a barrier or a join, so a thread that is preempted, interrupted or stopped by a page fault while it holds part of the work delays all the others, at every level that the delay overlaps. We present a wait-free method for direction-optimizing BFS, in which threads cooperate on a single search and every call returns after a bounded number of its own steps, whatever the other threads do.

Each level holds one bucket per thread. A thread expands its own bucket, then helps with what is left of the others, and orders its writes so that a thread that stops at any point leaves nothing that the others cannot redo. The threads agree on the direction of each level with a single compare-and-swap, bottom-up levels divide the vertices into blocks that are owned and helped like buckets, and one-byte codes give the bottom-up step an exact test of the current level, even for threads that have fallen behind. We prove the method correct and wait-free under the memory model of x86, and describe a variant whose memory depends on the largest levels of a search rather than on all of them.

On a 128-core, two-socket machine, against GAPBS, GBBS and PASGAL, our method is the fastest on four of five graphs, by 1.33 to 1.89$\times$ over the fastest of the three; on the road network, whose search has 6,262 levels, the memory optimized variant is 1.09$\times$ faster than PASGAL. On LiveJournal, when one of 16 threads stalls for 10~ms at the start of every level or round, GAPBS, GBBS and PASGAL take 2.8 to 9.0$\times$ as long as without stalls, while our method takes 16\% longer.
\end{abstract}
